\section{Introduction}\label{sec:intro}
Image segmentation is a fundamental computer vision task that partitions an image into meaningful, uniformly connected regions \cite{gonzalez2009digital}. It is a key preprocessing step for high-level computer vision applications, such as lane-tracking in autonomous vehicles \cite{janai2020computer}, object detection \cite{redmon2016you}, and medical imaging overlays \cite{pham2000current}. While simple color-based segmentation is common, modern real-time applications require spatial consistency. To achieve this, each pixel is mapped to a 5D feature manifold:
\begin{equation}
\mathbf{x} = [B, G, R, x, y]^T
\end{equation}
By clustering pixels within this combined color-spatial manifold, the algorithm naturally groups pixels that are both chromatically similar and geometrically close, preventing pixel-level noise \cite{achanta2012slic}.

Real-time video processing places strict constraints on execution latency. In computer vision, standard real-time throughput is governed by the human-facing display refresh rate of 60 frames per second (FPS) ($16.67\text{ ms}$ per frame) \cite{DHANANI20135}. Yet, for high-speed interactive systems, including driving simulators \cite{cruden2019simulator}, virtual reality headsets \cite{jerald2009scene}, and autonomous vehicle obstacle avoidance, the standard target is high-speed real-time processing operating at $120\text{ FPS}$ or more ($8.33\text{ ms}$ or less per frame) to prevent feedback loop delays and tracking oscillation \cite{cruden2019simulator, jerald2009scene}. Running heavy deep-learning segmentation models on embedded edge devices is often impossible under these extreme limits due to massive power consumption and high latency \cite{paszke2016enet}. Therefore, highly optimized, lightweight, Graphics Processing Unit (GPU)-accelerated classical clustering is a practical alternative for edge devices \cite{farivar2008parallel}.

Simultaneously, quantum-inspired computing offers a practical model for simulating quantum mechanics on classical hardware. True quantum computers are still in the Noisy Intermediate-Scale Quantum (NISQ) era \cite{preskill2018quantum}, characterized by small physical qubit registers, lack of error correction, and short coherence times ($T_1, T_2$) \cite{bharti2022noisy}. Physical hardware is highly sensitive to environmental noise, and transferring live video frames to a Quantum Processing Unit (QPU) introduces massive network latency, rendering genuine quantum hardware unusable for real-time video \cite{fraunhofer2025quantum}. To bypass these hardware constraints while leveraging quantum mathematical behaviors, quantum-inspired computing translates quantum mechanics such as state-overlap fidelity via the SWAP test \cite{buhrman2001quantum} (which measures $|\langle\psi\vert\phi\rangle|^2$) into deterministic equations that are parallelized directly inside custom GPU kernels. This allows for a direct comparison of the mathematical properties of classical and quantum clustering in real-time on standard local hardware.

\subsection{Problem Statement}
Designing a real-time, comparative 5D clustering framework introduces three major technical challenges:
\begin{enumerate}
    \item \textbf{Computational Complexity:} Evaluating 5D similarity for a standard Video Graphics Array (VGA) frame ($N = 307,200$ pixels) across $K$ clusters requires millions of calculations per frame. Doing this sequentially on a Central Processing Unit (CPU) is too slow for real-time feeds, resulting in massive lag.
    \item \textbf{NISQ Hardware Limitations:} While quantum algorithms can evaluate state overlaps efficiently, physical QPUs cannot operate within the sub-millisecond execution windows required for video pipelines due to hardware noise \cite{bharti2022noisy} and slow communication interfaces \cite{fraunhofer2025quantum}.
    \item \textbf{Lack of a Unified Benchmark:} To the best of the author's knowledge, there is no standard, open-source tool that allows for a direct and real-time comparison between classical Euclidean distance \cite{lloyd1982least} and quantum-inspired phase-space metrics under identical initialization and hardware constraints.
\end{enumerate}

\subsection{Objectives and Original Contributions}
This work develops and validates a high-throughput, GPU-accelerated C++ and Compute Unified Device Architecture (CUDA) framework for real-time 5D image segmentation. The original contributions of this work include:
\begin{enumerate}
    \item \textbf{An Optimized Quantum-Inspired Cosine Kernel:} Emulation of the quantum SWAP test circuit \cite{buhrman2001quantum} on the GPU by mapping pixel coordinates to phase states and calculating overlaps using direct hardware-accelerated cosine intrinsics (\texttt{\_\_cosf()}), leveraging Special Function Unit (SFU) hardware acceleration \cite{CudaGuide} to achieve high-throughput single-instruction execution.
    \item \textbf{A Modular Strategy Engine Architecture:} A fully decoupled C++ execution framework based on the Strategy design pattern \cite{gamma1995design}, allowing for seamless, real-time hot-swapping between the classical Euclidean and quantum-inspired GPU engines without resetting the video stream or losing state.
    \item \textbf{An Asynchronous Still-Frame Diagnostic Engine:} An on-demand diagnostic worker thread that captures a snapshot still-frame from the active feed to calculate heavy validation metrics (Within-Cluster Sum of Squares (WCSS) \cite{krzanowski1988criterion}, the Davies-Bouldin Index \cite{davies1979cluster}, and a Monte Carlo-approximated Silhouette score \cite{rousseeuw1987silhouettes}) without lagging the live video feed.
    \item \textbf{A Parallel Local-to-Global Reduction Loop:} A two-tier aggregation strategy that accumulates sub-totals in fast shared memory before updating global Video Random Access Memory (VRAM), reducing global memory traffic by a factor of $256\times$ \cite{CudaGuide}.
    \item \textbf{A Custom GPU-Native Preprocessor:} The \texttt{StridedDataPreprocessor}, which normalizes Blue-Green-Red (BGR) and spatial pixel coordinates and formats them into a contiguous Array of Structures (AoS) on the GPU, maximizing $L_1$ cache hit rates \cite{CudaGuide}.
\end{enumerate}

The rest of this paper is organized as follows: Section~\ref{sec:theory} reviews the theoretical foundations of 5D manifold scaling and quantum SWAP test emulation; Section~\ref{sec:architecture} presents the decoupled system architecture; Section~\ref{sec:implementation} details GPU implementation and optimizations; Section~\ref{sec:results} analyzes experimental performance and quality benchmarks; and Section~\ref{sec:conclusion} concludes the paper.
